% !TeX program = xelatex
\documentclass[11pt,a4paper]{article}
\usepackage[margin=2.1cm]{geometry}
% XeLaTeX/LuaLaTeX dùng font hệ thống; pdfLaTeX dùng mã T5 cho tiếng Việt.
\usepackage{iftex}
\ifPDFTeX
  \usepackage[utf8]{inputenc}
  \usepackage[T5]{fontenc}
  \usepackage[vietnamese]{babel}
\else
  \usepackage{fontspec}
  \IfFontExistsTF{Times New Roman}
    {\setmainfont{Times New Roman}}
    {\setmainfont{texgyretermes-regular.otf}[
      BoldFont=texgyretermes-bold.otf,
      ItalicFont=texgyretermes-italic.otf,
      BoldItalicFont=texgyretermes-bolditalic.otf]}
  \IfFontExistsTF{Arial}
    {\setsansfont{Arial}}
    {\setsansfont{Latin Modern Sans}}
  \IfFontExistsTF{Courier New}
    {\setmonofont{Courier New}[Scale=0.85]}
    {\setmonofont{Latin Modern Mono}[Scale=0.85]}
\fi
\usepackage{graphicx,booktabs,tabularx,array,enumitem,xcolor}
\usepackage{tikz}
\usetikzlibrary{arrows.meta,positioning}
\usepackage{hyperref}
\hypersetup{colorlinks=true,urlcolor=blue!55!black,linkcolor=black}
\usepackage{fancyhdr}
\pagestyle{fancy}
\fancyhf{}
\fancyhead[L]{\small Tích hợp AI vào hệ thống VMS}
\fancyhead[R]{\small Module FACE}
\fancyfoot[C]{\thepage}
\setlength{\headheight}{14pt}
\setlength{\parindent}{0pt}
\setlength{\parskip}{5pt}
\setlist[itemize]{leftmargin=1.4em,itemsep=2pt,topsep=3pt}
\renewcommand{\figurename}{Hình}
\renewcommand{\tablename}{Bảng}
\renewcommand{\contentsname}{Mục lục}
\renewcommand{\refname}{Tài liệu tham khảo}
\newcolumntype{Y}{>{\raggedright\arraybackslash}X}
% Điền ba trường dưới đây, không cần sửa phần bố cục.
\newcommand{\TenBaoCao}{}
\newcommand{\HoTen}{}
\newcommand{\Github}{}
\newcommand{\truong}[1]{\if\relax\detokenize{#1}\relax\rule{0.70\linewidth}{0.4pt}\else#1\fi}
\newcommand{\anhminhchung}[2]{%
  \IfFileExists{#1}{\includegraphics[width=\linewidth,height=4.3cm,keepaspectratio]{#1}}{%
  \fbox{\parbox[c][3.4cm][c]{0.94\linewidth}{\centering\textcolor{gray}{#2}\\[4pt]
  \small\texttt{\detokenize{#1}}}}}}
\begin{document}
% Phần đầu gọn, nội dung bắt đầu ngay trên cùng trang.
\thispagestyle{plain}
\begin{center}
{\large\bfseries BÁO CÁO THỰC TẬP TUẦN 2\par}
\vspace{4pt}
{\LARGE\bfseries
\if\relax\expandafter\detokenize\expandafter{\TenBaoCao}\relax
  \rule{0.72\linewidth}{0.4pt}% Để trống tên báo cáo để điền sau.
\else
  \TenBaoCao
\fi\par}
\end{center}
\vspace{3pt}
\begingroup
\setlength{\parskip}{3pt}
\textbf{Họ và tên:} \expandafter\truong\expandafter{\HoTen}\par
\textbf{Chức năng:} \texttt{FACE} -- Nhận diện khuôn mặt\par
\textbf{Môi trường:} AIOC/VMS staging\par
\textbf{GitHub:} \expandafter\truong\expandafter{\Github}\par
\endgroup
\vspace{2pt}

\section{Mục tiêu và yêu cầu bài toán}
Project xây dựng module nhận diện khuôn mặt từ camera RTSP và tích hợp vào VMS staging qua MQTT. Kết quả cần có khung phát hiện trên video, sự kiện nhận diện và ảnh snapshot trên Web VMS.

Phạm vi chức năng gồm phát hiện khuôn mặt, theo dõi đối tượng, so khớp với danh sách mẫu tối thiểu 5 nhân viên và xác định hướng IN/OUT theo vùng hoặc vạch cấu hình. Với người nhận diện được, sự kiện gửi mã nhân viên, tên và điểm so khớp. Với người lạ, sự kiện không gửi mã nhân viên.

Các yêu cầu tích hợp chính là nhận cấu hình camera/vùng từ VMS, chỉ xử lý camera bật FACE, cập nhật cấu hình khi có thay đổi, upload ảnh qua Storage API và publish đúng topic. Sự kiện chỉ được tạo từ detection thật, có kiểm soát lặp bằng cooldown. Database được triển khai bằng Docker Compose; thông tin kết nối đặt trong \texttt{.env}.

\section{Kiến trúc và luồng xử lý}
\begin{figure}[ht]
\centering
\begin{tikzpicture}[node distance=7mm and 5mm,
 box/.style={draw,rounded corners,align=center,font=\small,minimum height=9mm,text width=3cm},
 arr/.style={-{Latex[length=2mm]},thick}]
\node[box] (vms) {VMS: camera, vùng/vạch};
\node[box,right=of vms] (mqtt) {VMS client\\nhận cấu hình MQTT};
\node[box,right=of mqtt] (read) {Đọc frame RTSP\\giữ frame mới nhất};
\node[box,below=of read] (detect) {YOLO11n TensorRT\\phát hiện khuôn mặt};
\node[box,left=of detect] (track) {ByteTrack\\gán track ID};
\node[box,left=of track] (id) {Extractor + pgvector\\so khớp danh tính};
\node[box,below=of id] (event) {Event processor\\polygon/line, IN/OUT};
\node[box,right=of event] (send) {Upload snapshot\\publish FACE và bbox};
\node[box,right=of send] (web) {Backend / Web VMS\\lưu và hiển thị};
\draw[arr] (vms)--(mqtt); \draw[arr] (mqtt)--(read);
\draw[arr] (read)--(detect); \draw[arr] (detect)--(track);
\draw[arr] (track)--(id); \draw[arr] (id)--(event);
\draw[arr] (event)--(send); \draw[arr] (send)--(web);
\end{tikzpicture}
\caption{Luồng xử lý của module FACE.}
\end{figure}

Mỗi camera có reader, tracker và cache identity riêng. Detector, extractor và kết nối database được dùng chung; các camera được xử lý lần lượt. Reader chỉ giữ frame mới nhất để hạn chế tích lũy ảnh cũ.

Database PostgreSQL của module AI lưu nhân viên và embedding phục vụ so khớp. Backend VMS tiếp nhận và lưu sự kiện qua MQTT. Hai thành phần này có vai trò khác nhau; danh sách mẫu của module AI chưa tự đồng bộ với danh sách nhân viên trên VMS.

\section{Phương pháp triển khai}
\subsection{Phát hiện và theo dõi khuôn mặt}
Detector sử dụng checkpoint YOLO11n đã huấn luyện cho một lớp khuôn mặt. Ảnh được đưa vào detector ở kích thước 640, với ngưỡng confidence 0,5. TensorRT FP16 được lựa chọn để cân bằng tốc độ, kích thước engine và chất lượng phát hiện.

ByteTrack gán ID cho các khuôn mặt qua nhiều frame. Kết quả tracker được ghép lại với detection của frame hiện tại. Chỉ bbox từ detection thật được sử dụng cho nhận diện và sự kiện; bbox dự đoán không tạo sự kiện. Mỗi lần camera nối lại hoặc tracker được tạo lại, lịch sử liên quan được xóa.

\subsection{Đăng ký và nhận diện nhân viên}
File \texttt{configs/employees.csv} chứa mã, tên và đường dẫn ảnh của 5 nhân viên. Script đăng ký kiểm tra mỗi ảnh có đúng một khuôn mặt, cắt ảnh mặt và tạo embedding bằng cùng extractor dùng trong pipeline. Ảnh không hợp lệ hoặc mặt quá nhỏ được bỏ qua và ghi log. Chạy lại cùng mã nhân viên sẽ cập nhật tên và embedding.

Extractor ONNX nhận ảnh mặt RGB kích thước $112\times112$, chuẩn hóa theo mean/std của model, tạo vector 512 chiều và chuẩn hóa L2. Vector được lưu vào cột \texttt{embedding} kiểu \texttt{vector(512)} trong bảng \texttt{employees}.

Khi nhận diện, pgvector tìm nhân viên có cosine similarity cao nhất:
\[
s(\mathbf{q},\mathbf{e})=
\frac{\mathbf{q}\cdot\mathbf{e}}{\|\mathbf{q}\|_2\|\mathbf{e}\|_2}.
\]
Nếu điểm cao nhất đạt ngưỡng 0,4 thì kết quả được xem là khớp; ngược lại là người lạ. Ngưỡng này là cấu hình thử nghiệm, chưa được hiệu chỉnh bằng một bộ kiểm thử nhận diện riêng. Hiện mỗi người có một ảnh mẫu và ảnh mặt chưa được căn chỉnh theo landmarks.

\subsection{Xử lý polygon, line crossing và IN/OUT}
Điểm đại diện là tâm bbox khuôn mặt, không phải vị trí chân người. Với polygon, ngoài vào trong là IN, trong ra ngoài là OUT. Với line crossing, \texttt{direction\_point} từ VMS xác định phía IN. Đối tượng phải cắt đoạn vạch hữu hạn; cắt đường kéo dài ngoài hai đầu vạch không được tính.

Track mới xuất hiện chỉ ghi nhận vị trí ban đầu. Các frame tiếp theo mới được dùng để kiểm tra chuyển phía. Dải đệm biên 0,003 giúp giảm rung bbox; cooldown 1 giây áp dụng theo camera, vùng và track. Track không còn detection quá 2 giây bị loại khỏi lịch sử. Các tham số đều có thể cấu hình.

Cơ chế này hạn chế sự kiện lặp khi cùng track đứng tại biên. Tuy nhiên, nếu tracker đổi ID cho cùng một người thì vẫn có thể tạo lượt mới. Polygon cũng có thể bỏ sót khi đối tượng đi xuyên toàn bộ vùng giữa hai frame mà cả hai vị trí quan sát đều ở ngoài.

\subsection{Cache identity}
Mỗi track được nhận diện ngay khi xuất hiện, sau đó dùng lại identity trong 0,5 giây trước khi làm mới. Cache lưu kết quả danh tính, không lưu ảnh hay embedding. Cả kết quả người quen và người lạ đều được cập nhật định kỳ.

Khi phát sinh IN/OUT, identity được làm mới trên chính frame tạo sự kiện trước khi publish. Nếu vừa nhận diện tại frame đó thì không extract lần thứ hai. Extractor vẫn xử lý từng mặt, không dùng batch. Cache được xóa khi track hết hạn hoặc tracker reset. Chưa có phép đo riêng để định lượng mức giảm độ trễ của cache.

\section{Tích hợp VMS và triển khai}
\subsection{Giao tiếp MQTT và snapshot}
VMS client kết nối MQTT qua WebSocket bảo mật, dùng client ID riêng, tự kết nối và subscribe lại khi mất kết nối. Danh sách camera được xin theo zone 51, organization 103.

\begin{table}[ht]
\centering\small
\begin{tabularx}{\linewidth}{@{}Y c@{}}
\toprule
Topic và chức năng & QoS\\
\midrule
\texttt{smart\_vms/cameras/getdevice}: xin camera & 1\\
\texttt{smart\_vms/cameras/zone/51}: nhận camera & 1\\
\texttt{smart\_vms/cameras/+/ai\_zones/FACE}: nhận vùng/vạch & 1\\
\texttt{smart\_vms/ai\_events/FACE}: gửi sự kiện & 1\\
\texttt{smart\_vms/ai/bbox/\{camera\_code\}}: gửi bbox & 0\\
\bottomrule
\end{tabularx}
\caption{Các topic được module sử dụng.}
\end{table}

Topic \texttt{ai\_zones/FACE} chứa cả polygon và line crossing theo hợp đồng VMS. Payload rỗng được xử lý như xóa vùng. Camera tắt FACE hoặc bị xóa được dừng xử lý ở vòng cập nhật tiếp theo.

Khi có sự kiện, ảnh JPEG được upload qua Storage API. Giá trị \texttt{object\_key} trả về được đặt vào \texttt{snapshot\_url}; ảnh không được gửi base64 qua MQTT. Payload FACE gồm module, camera, thời gian có múi giờ \texttt{+07:00}, hướng IN/OUT và thông tin nhận diện. Trường \texttt{user\_code} chỉ được gửi khi nhận diện thành công.

BBox được chuẩn hóa về khoảng 0--1 và gửi theo frame. Field \texttt{color} mang màu xanh khi khớp và đỏ khi là người lạ. Việc hiển thị màu phụ thuộc frontend đọc field này. Bbox và sự kiện FACE là hai luồng riêng; gửi bbox không tự tạo lịch sử nhận diện.

\subsection{Docker và đối chiếu yêu cầu}
Docker Compose triển khai PostgreSQL 16 có pgvector, sử dụng schema trong \texttt{sql/init.sql} và volume lưu dữ liệu. Trong phạm vi hiện tại, database chạy container; chương trình AI được chạy bằng môi trường Python trên máy.

\begin{table}[ht]
\small
\begin{tabularx}{\linewidth}{@{}l Y Y@{}}
\toprule
Mã & Cách thực hiện & Bằng chứng và giới hạn\\
\midrule
C01 & MQTT WSS, reconnect, subscribe lại & Đã triển khai; có kiểm thử callback offline.\\
C02 & Lọc camera bật FACE, cập nhật danh sách & Đã triển khai; cần bổ sung minh chứng thao tác bật/tắt trên VMS.\\
C03 & Cập nhật vùng/vạch, xử lý rỗng & Đã triển khai; log có camera nhận được một vùng FACE.\\
C04 & Xử lý luồng RTSP & Demo đã chạy trên luồng video thực tế.\\
C05 & Upload JPEG, gửi object key & Log có object key snapshot cho 6 sự kiện FACE.\\
C06 & Topic FACE, thời gian có timezone & Log có publish thành công; cần ảnh xác nhận lịch sử trên Web VMS.\\
C07 & Bbox chuẩn hóa, QoS 0 & Đã có log publish; cần ảnh bbox trên Web VMS.\\
C08 & Chuyển phía, cooldown theo track & Có kiểm thử logic; chưa thống kê trùng/bỏ sót trên video có nhãn.\\
C09 & Chỉ detection thật & Tracker ghép detection hiện tại; có kiểm thử bỏ bbox dự đoán.\\
\bottomrule
\end{tabularx}
\caption{Đối chiếu triển khai với yêu cầu của đề bài.}
\end{table}

\section{Thực nghiệm và đánh giá}
\subsection{Đánh giá detector và TensorRT}
Thử nghiệm chạy trên NVIDIA GeForce RTX 4050 Laptop 6 GB, PyTorch 2.13.0 với CUDA 12.6, Ultralytics 8.4.123 và TensorRT 10.16.1.11. Các biến thể dùng cùng checkpoint YOLO11n, input $640\times640$, batch 1.

Timing sử dụng 24 ảnh test nạp sẵn, warmup 20 lần và đo 3 vòng, mỗi vòng 100 lần. CUDA được đồng bộ trước và sau phép đo. Latency predict bao gồm tiền xử lý, forward, hậu xử lý và lấy bbox/score về CPU; không bao gồm đọc ảnh, load model, extractor, database hoặc MQTT.

INT8 được calibration bằng 500 ảnh train. Độ chính xác detector được đánh giá trên 164 ảnh test với 263 khuôn mặt được gán nhãn, tách khỏi calibration.

\begin{table}[ht]
\centering\small
\begin{tabular}{@{}lrrrrr@{}}
\toprule
Model & ms/ảnh & p95 (ms) & FPS & mAP50 & mAP50--95\\
\midrule
PyTorch FP32 & 14,54 & 16,11 & 68,76 & 95,32\% & 67,27\%\\
TRT FP32 mặc định & 7,45 & 8,42 & 134,29 & 95,32\% & 67,28\%\\
\textbf{TRT FP16} & \textbf{7,24} & \textbf{8,68} & \textbf{138,05} & \textbf{95,33\%} & \textbf{67,35\%}\\
TRT INT8 & 6,93 & 8,17 & 144,21 & 94,25\% & 66,84\%\\
\bottomrule
\end{tabular}
\caption{Kết quả detector; FPS được tính từ latency predict trung bình.}
\end{table}

TensorRT FP16 giảm latency predict khoảng 50,2\% so với PyTorch FP32 và đạt tốc độ khoảng 2,01 lần. Engine FP16 có kích thước 7,97 MB, nhỏ hơn engine TRT FP32 12,33 MB. Chênh lệch mAP giữa FP16 và baseline rất nhỏ trên bộ test này.

INT8 nhanh nhất nhưng chỉ giảm thêm khoảng 0,31 ms/ảnh so với FP16; mAP50 thấp hơn baseline khoảng 1,08 điểm phần trăm. Vì vậy, FP16 được lựa chọn để cân bằng tốc độ, kích thước engine và chất lượng phát hiện. Kích thước file không được xem là phép đo VRAM khi chạy.

Recall tổng hợp của FP16 là 89,46\%, thấp hơn baseline 92,63\%, trong khi Precision tăng. Các chỉ số này được Ultralytics lấy ở điểm F1 tốt nhất của từng biến thể, nên cần đánh giá thêm tại cùng confidence triển khai để kết luận mức bỏ sót. Không suy ra độ chính xác nhận diện danh tính từ mAP detector.

\begin{figure}[ht]
\centering
\IfFileExists{images/latency_comparison.png}
  {\includegraphics[width=\linewidth]{images/latency_comparison.png}}
  {\anhminhchung{images/latency_comparison.png}{Bổ sung biểu đồ benchmark detector}}
\caption{So sánh latency và throughput của riêng detector trên GPU thử nghiệm.}
\end{figure}

\subsection{Đánh giá nhận diện danh tính}
Model trích xuất đặc trưng hiện còn ở giai đoạn thử nghiệm, chưa được huấn luyện và tinh chỉnh đầy đủ cho điều kiện camera thực tế. Khi chạy thử, khả năng nhận diện chưa ổn định; một số người đã đăng ký vẫn bị phân loại là người lạ. Danh sách mẫu hiện có 5 nhân viên, mỗi người một ảnh, nên chưa bao quát được các góc mặt và điều kiện ánh sáng khác nhau.

Nhận định ban đầu là model và dữ liệu mẫu cần được cải thiện thêm. Tuy nhiên, chưa có thực nghiệm để xác định mức ảnh hưởng riêng của quá trình huấn luyện, chất lượng ảnh và ngưỡng so khớp. Phần nhận diện danh tính hiện mới được kiểm tra định tính qua chạy thử, chưa đo tỷ lệ nhận đúng, nhận nhầm hoặc từ chối người đã đăng ký. Vì vậy, báo cáo chưa đưa ra số liệu độ chính xác cho phần này. Hướng tiếp theo là tiếp tục huấn luyện, bổ sung ảnh mẫu và đánh giá trên tập dữ liệu độc lập.

\subsection{Sự kiện và bằng chứng tích hợp}
Log ngày 30/10/2026, khoảng 22:38:22 đến 22:39:01, ghi nhận 6 lần publish FACE của \texttt{camera\_face\_01}. Các sự kiện có hướng IN, trạng thái người lạ và object key snapshot. Kết quả publish là thành công ở mức broker ACK, chưa phải xác nhận backend đã lưu vào database. Camera thứ hai trong đoạn log chưa có vùng FACE nên không tạo IN/OUT.

Hai hình dưới dành cho bằng chứng Web VMS. Các ảnh này chưa có trong bộ tài liệu báo cáo hiện tại và cần bổ sung trước khi nộp. Chưa có thống kê ground truth để kết luận số lượt IN/OUT đúng, trùng hoặc bỏ sót.

\begin{figure}[ht]
\centering
\anhminhchung{images/vms_bbox.png}{Bổ sung ảnh bbox FACE đang hiển thị trên Web VMS}
\caption{Minh chứng bbox trên trang giám sát trực tiếp, cần bổ sung ảnh.}
\end{figure}
\begin{figure}[ht]
\centering
\anhminhchung{images/vms_event.png}{Bổ sung ảnh lịch sử FACE và snapshot trên Web VMS}
\caption{Minh chứng sự kiện và ảnh snapshot trên Web VMS, cần bổ sung ảnh.}
\end{figure}

\section{Kết luận và hướng phát triển}
Với project này, pipeline từ đọc RTSP, phát hiện, tracking, tạo embedding, so khớp database đến xử lý IN/OUT và publish MQTT đã được xây dựng. Script enroll hỗ trợ cập nhật dữ liệu từ CSV. TensorRT FP16 được lựa chọn sau benchmark; cache identity giảm số lần gọi extractor và search, nhưng hiệu quả độ trễ toàn pipeline chưa được đo.

Các hạn chế còn lại gồm ít ảnh mẫu, chưa căn chỉnh khuôn mặt, ngưỡng so khớp chưa hiệu chỉnh, khả năng đổi track ID và upload snapshot đang chặn vòng xử lý. Bằng chứng lưu/hiển thị sự kiện trên Web VMS và bảng đánh giá nhận diện cần được bổ sung.

Hướng phát triển ưu tiên là tiếp tục huấn luyện và tinh chỉnh model trích xuất đặc trưng với dữ liệu phù hợp camera thực tế, nhằm cải thiện tình trạng người đã đăng ký vẫn bị nhận là người lạ. Đồng thời, cần bổ sung nhiều ảnh mẫu cho mỗi nhân viên ở các góc mặt và điều kiện ánh sáng khác nhau, xem xét căn chỉnh khuôn mặt trước khi tạo embedding.

Sau khi cải thiện model và dữ liệu, cần xây dựng tập kiểm thử độc lập gồm người đã đăng ký và người lạ, đo tỷ lệ nhận đúng, nhận nhầm và từ chối người đã đăng ký để lựa chọn ngưỡng cosine phù hợp. Các hướng tối ưu hệ thống gồm kiểm tra bỏ sót ở confidence triển khai, đo thời gian từng bước của toàn pipeline và tách upload/publish sự kiện sang worker riêng.

\clearpage
\section*{Phụ lục: thao tác chạy}
\addcontentsline{toc}{section}{Phụ lục: thao tác chạy}
Từ thư mục project, kích hoạt môi trường Python có CUDA, Ultralytics, TensorRT và các thư viện của hệ thống, sau đó chạy:
\begin{quote}\small\ttfamily
docker compose up -d db\\
python -m scripts.enroll\_employees -{}-image-root D:/private/task2\_atin\\
python main.py
\end{quote}
Lệnh enroll trên sử dụng thư mục ảnh hiện có của project cũ. Nếu đã chuyển ảnh vào thư mục project mới thì cần điều chỉnh \texttt{-{}-image-root}. Thông tin tài khoản không ghi trong báo cáo; các biến kết nối được cấu hình qua \texttt{.env}.

Notebook thử nghiệm: \texttt{experiments/detector\_tensorrt/}\\
\texttt{detector\_quantization\_benchmark.ipynb}.\\
Cấu hình chạy: \texttt{configs/config.yml}.\\
Video demo: \rule{0.66\linewidth}{0.4pt}

\begin{thebibliography}{9}\small
\bibitem{assignment} Tài liệu nội bộ, \emph{Bài tập thực tập tuần 2: Tích hợp AI vào hệ thống VMS}, ngày giao 26/10/2026.
\bibitem{mqtt} Tài liệu nội bộ, \emph{AIOC Staging -- AI Modules \& MQTT Contract}, các mục 4, 5.2, 5.8, 6 và 9.
\bibitem{ultralytics} Ultralytics, \emph{TensorRT Export}. \url{https://docs.ultralytics.com/integrations/tensorrt/}
\bibitem{pgvector} pgvector, \emph{Vector similarity search for Postgres}. \url{https://github.com/pgvector/pgvector}
\end{thebibliography}
\end{document}
